\uuid{pC8p}
\exo7id{6732}
\titre{exo7 6732}
\auteur{queffelec}
\organisation{exo7}
\datecreate{2011-10-16}
\isIndication{false}
\isCorrection{false}
\chapitre{Autre}
\sousChapitre{Autre}
\cpge{mpsi-mp2i.nombres-complexes.a, pcsi.nombres-complexes.a, ptsi.nombres-complexes.a, mpsi-mp2i.nombres-complexes.h, pcsi.nombres-complexes.h, ptsi.nombres-complexes.h, ptsi.geometrie-plan.d, bcpst1.nombres-complexes.a, bcpst1.geometrie.c}
\module{Analyse complexe}
\niveau{L3}
\difficulte{}

\contenu{
\texte{
Calculer $\int_C(z^2+3z)\ dz$ le long des chemins suivants :
}
\begin{enumerate}
    \item \question{le cercle $\vert z\vert=2$ du point (2,0) au point (0,2)}
    \item \question{le segment de droite joignant les points (2,0) et (0,2)}
    \item \question{le contour polygonal formé par les segments de droite joignant
(2,0) à (2,2) et (2,2) à (0,2)}
\end{enumerate}
}
